\documentclass[12pt,a4paper]{article}

\usepackage[utf8]{inputenc}
\usepackage[T1]{fontenc}
\usepackage[spanish]{babel}
\usepackage{geometry}
\usepackage{longtable}
\usepackage{xcolor}
\usepackage{hyperref}
\usepackage{listings}
\usepackage{tikz}

\usetikzlibrary{arrows.meta,positioning,babel}

\geometry{margin=2.5cm}

\lstdefinestyle{code}{
    basicstyle=\ttfamily\small,
    breaklines=true,
    frame=single,
    columns=fullflexible,
    showstringspaces=false
}

\title{\textbf{Proyecto 3EnRaya}\\Informe tecnico de Teleinformatica}
\author{
    Estudiante: \underline{\hspace{8cm}}\\
    Codigo: \underline{\hspace{8cm}}\\
    Docente: Luis Marcel Barrero Mendizabal\\
    Materia: Teleinformatica
}
\date{Mayo de 2026}

\begin{document}

\begin{titlepage}
    \centering
    {\Large Universidad \underline{\hspace{8cm}}\par}
    \vspace{1.5cm}
    {\Huge \textbf{Proyecto 3EnRaya}\par}
    \vspace{0.5cm}
    {\Large Informe tecnico de Teleinformatica\par}
    \vspace{2cm}
    \begin{flushleft}
    \textbf{Estudiante:} \underline{\hspace{9cm}}\\[0.4cm]
    \textbf{Codigo:} \underline{\hspace{9.8cm}}\\[0.4cm]
    \textbf{Docente:} Luis Marcel Barrero Mendizabal\\[0.4cm]
    \textbf{Materia:} Teleinformatica\\[0.4cm]
    \textbf{Fecha:} Mayo de 2026
    \end{flushleft}
    \vfill
    {\large Linux Mint - Lenguaje C - Sockets TCP\par}
\end{titlepage}

\tableofcontents
\newpage

\section{Introduccion}

3EnRaya es un sistema cliente-servidor desarrollado en lenguaje C. El servidor
central recibe conexiones TCP, registra jugadores, mantiene una cola de espera
y administra completamente una partida de tres en raya 1 vs 1.

El juego no se ejecuta entre pares directamente. Los clientes solo envian
comandos al servidor y reciben respuestas. El servidor valida turnos,
movimientos, tablero, ganador, empate y marcador.

\section{Objetivo}

Implementar un servicio de red con sockets TCP donde varios jugadores puedan
conectarse a un servidor central y participar en partidas de tres en raya.

\section{Arquitectura}

\begin{figure}[h]
\centering
\begin{tikzpicture}[
    box/.style={draw, rounded corners=4pt, very thick, align=center, minimum width=3cm, minimum height=1.2cm, fill=blue!8},
    server/.style={draw, rounded corners=4pt, very thick, align=center, minimum width=4cm, minimum height=1.5cm, fill=green!10},
    arrow/.style={{Stealth[length=3mm]}-{Stealth[length=3mm]}, very thick}
]
\node[server] (s) {\textbf{Servidor central}\\TCP :5000\\select + juego};
\node[box, left=3.4cm of s] (a) {\textbf{Cliente A}\\Jugador X};
\node[box, right=3.4cm of s] (b) {\textbf{Cliente B}\\Jugador O};
\node[box, below=2.2cm of s] (n) {\textbf{Cliente N}\\Lobby};
\draw[arrow] (a) -- node[above]{TCP} (s);
\draw[arrow] (s) -- node[above]{TCP} (b);
\draw[arrow] (n) -- node[right]{TCP} (s);
\end{tikzpicture}
\caption{Arquitectura cliente-servidor.}
\end{figure}

\section{Makefile}

El proyecto usa un Makefile inspirado en el formato de clase. Sus variables
principales son:

\begin{lstlisting}[style=code]
CC = gcc
CFLAGS = -std=c11 -Wall -Wextra -Wrestrict -pedantic -O2
EXT = .c
SRCDIR = src
OBJDIR = obj
BINDIR = bin
\end{lstlisting}

A diferencia del template de un solo ejecutable, este proyecto genera dos:

\begin{lstlisting}[style=code]
bin/3enraya_server
bin/3enraya_client
\end{lstlisting}

\section{Protocolo}

El protocolo es de texto y cada mensaje termina en salto de linea. Los comandos
principales son:

\begin{longtable}{|p{4cm}|p{9cm}|}
\hline
\textbf{Comando} & \textbf{Descripcion} \\
\hline
\texttt{HELLO <nombre>} & Registra al jugador. \\
\hline
\texttt{QUEUE} & Entra a cola de espera. \\
\hline
\texttt{PLAYERS} & Lista participantes conectados. \\
\hline
\texttt{SCORE} & Muestra marcador. \\
\hline
\texttt{BOARD} & Muestra tablero y turno. \\
\hline
\texttt{MOVE <1-9>} & Realiza una jugada. \\
\hline
\texttt{QUIT} & Cierra la conexion. \\
\hline
\end{longtable}

\section{Estructuras de datos}

\subsection{\texttt{Player}}

Representa un jugador conectado. Contiene socket, direccion IP, nombre, buffer
de recepcion, estado de cola, partida actual, simbolo y marcador.

\subsection{\texttt{Game}}

Representa una partida activa. Contiene identificador, jugador X, jugador O,
tablero y turno actual.

\section{Estados del jugador}

\begin{lstlisting}[style=code]
DESCONECTADO -> CONECTADO -> LOBBY -> EN_COLA -> EN_PARTIDA -> LOBBY
\end{lstlisting}

Desde cualquier estado conectado, el jugador puede pasar a \texttt{DESCONECTADO}
por \texttt{QUIT}, error o desconexion.

\section{Compilacion y ejecucion}

\begin{lstlisting}[style=code]
make
./bin/3enraya_server --host 127.0.0.1 --port 5000
./bin/3enraya_client --host 127.0.0.1 --port 5000 --name Ramiro
./bin/3enraya_client --host 127.0.0.1 --port 5000 --name Paulo
\end{lstlisting}

\section{Conclusiones}

\begin{itemize}
    \item El servidor central controla completamente el juego.
    \item \texttt{select} permite atender varios clientes sin procesos ni hilos.
    \item El Makefile separa compilacion de servidor y cliente usando objetos.
    \item El protocolo de texto facilita la defensa y la depuracion.
    \item El proyecto es compatible con Visual Studio Code y extensiones de C/C++.
\end{itemize}

\end{document}
